\documentclass[11pt]{article}
\usepackage[margin=1in]{geometry}
\usepackage{booktabs}
\usepackage{longtable}
\usepackage{hyperref}
\usepackage{siunitx}
\usepackage{fancyhdr}
\usepackage{lastpage}
\pagestyle{fancy}
\fancyhf{}
\rhead{Tri-State Systems Manager}
\lhead{Engineering Evidence Package}
\cfoot{Page \thepage\ of \pageref{LastPage}}

\title{Community River Valley Engineering Evidence Package}
\author{Tri-State Systems Manager}
\date{Generated from a versioned evidence manifest}

\begin{document}
\maketitle

\section*{Authority and review boundary}
This document is a reproducible engineering evidence package. It is not a professional engineering certification, permit, flood-insurance determination, levee accreditation, emergency order, or agency acceptance. Human engineers and the responsible agencies retain final authority.

\section{Project identity}
Project ID: \texttt{\detokenize{<PROJECT_ID>}}\\
Model version: \texttt{\detokenize{<MODEL_VERSION>}}\\
Equation set: \texttt{\detokenize{<EQUATION_SET>}}\\
Numerical tolerance: \texttt{\detokenize{<NUMERICAL_TOLERANCE>}}\\
Generated: \texttt{\detokenize{<GENERATED_AT>}}

\section{Hydrology and hydraulics}
Record the observed station IDs, observation timestamps, parameter codes, units, vertical datum, terrain source, boundary conditions, Manning roughness assumptions, mesh/breakline metadata, structure representation, water-surface elevations, depths, velocities, shear stresses, no-rise comparison and uncertainty/sensitivity results. Every value must reference an evidence identifier.

\section{Subsurface-to-finished-grade section}
Record survey control, mapped-soil context, site investigation, groundwater/pore pressure, weak or organic strata, foundation preparation, geosynthetics/drainage, qualified fill, compacted lifts, pavement/aggregate, drainage structures, design crest/freeboard, erosion/scour protection and instrumentation. Generalized soil mapping is context and does not replace site investigation.

\section{Dredged-material qualification}
Record dredging source and campaign, sample chain of custody, gradation, moisture, Atterberg limits, organic content, density, compaction, shear strength, consolidation/compressibility, hydraulic conductivity, contaminant/leachability testing, dewatering/processing, blending and acceptance. Structural placement remains prohibited until the applicable engineering and environmental gates are satisfied.

\section{Regulatory and funding gates}
Record applicable Indiana DNR, FEMA, Clean Water Act, navigation, wetlands/species/cultural-resource, right-of-way, utility, stormwater, professional survey/engineering and grant-program requirements. Each gate must contain its authoritative source, current verification date and unresolved actions.

\section{Cryptographic provenance}
Input SHA-256: \texttt{\detokenize{<INPUT_HASH>}}\\
Output SHA-256: \texttt{\detokenize{<OUTPUT_HASH>}}\\
Manifest SHA-256: \texttt{\detokenize{<MANIFEST_HASH>}}

\section{Review state}
Software-generated packages use \texttt{ENGINEERING\_REVIEW\_READY} as the highest automated state. Any \texttt{ENGINEER\_ACCEPTED} or \texttt{AGENCY\_ACCEPTED} state must be recorded from an authorized human reviewer and never inferred by software.

\end{document}
